\documentclass{article}
\usepackage{amsmath, amssymb, amsthm}
\usepackage{hyperref}
\usepackage{geometry}
\geometry{margin=1in}

\title{Erd\H{o}s Problem \#600}
\date{Last updated: 2025-08-31}
\author{Source: erdosproblems.com}

\begin{document}
\maketitle

\noindent\textbf{Status:} OPEN \quad \textbf{No prize}

\noindent\textbf{Tags:} graph theory

\medskip

\noindent\textbf{Problem Statement:}

Let $e(n,r)$ be minimal such that every graph on $n$ vertices with at least $e(n,r)$ edges, each edge contained in at least one triangle, must have an edge contained in at least $r$ triangles. Let $r\geq 2$. Is it true that\[e(n,r+1)-e(n,r)\to \infty\]as $n\to \infty$? Is it true that\[\frac{e(n,r+1)}{e(n,r)}\to 1\]as $n\to \infty$?




Ruzsa and Szemer\'{e}di \cite{RuSz78} proved that $e(n,r)=o(n^2)$ for any fixed $r$. See also [80] .




References

[RuSz78] Ruzsa, I. Z. and Szemer\'{e}di, E., Triple systems with no six points carrying three triangles . Combinatorics (Proc. Fifth Hungarian Colloq., Keszthely, 1976), Vol. II (1978), 939-945.

\noindent\textbf{Related OEIS sequences:} possible


\bigskip
\noindent\small{Source: \url{https://www.erdosproblems.com/600}}
\end{document}
